% ARQUIVO ÚNICO: tabelas e referências já estão dentro deste arquivo.
% No Overleaf, basta este arquivo (e, se quiser, figuras/pertinencias.png).
\begin{filecontents*}[overwrite]{refs.bib}
% Confira cada entrada no PDF/site do artigo antes de submeter. Entradas com TODO precisam ser completadas.

@inproceedings{chen2019,
  author    = {Chen, Tongfei and Navr{\'a}til, Ji{\v{r}}{\'\i} and Iyengar, Vijay and Shanmugam, Karthikeyan},
  title     = {Confidence Scoring Using Whitebox Meta-models with Linear Classifier Probes},
  booktitle = {Proceedings of the 22nd International Conference on Artificial Intelligence and Statistics (AISTATS)},
  series    = {Proceedings of Machine Learning Research},
  volume    = {89},
  pages     = {1467--1475},
  publisher = {PMLR},
  year      = {2019}
}

@inproceedings{corbiere2019,
  author    = {Corbi{\`e}re, Charles and Thome, Nicolas and Bar-Hen, Avner and Cord, Matthieu and P{\'e}rez, Patrick},
  title     = {Addressing Failure Prediction by Learning Model Confidence},
  booktitle = {Advances in Neural Information Processing Systems (NeurIPS)},
  year      = {2019}
}

@inproceedings{hendrycks2017,
  author    = {Hendrycks, Dan and Gimpel, Kevin},
  title     = {A Baseline for Detecting Misclassified and Out-of-Distribution Examples in Neural Networks},
  booktitle = {International Conference on Learning Representations (ICLR)},
  year      = {2017}
}

@inproceedings{gal2016,
  author    = {Gal, Yarin and Ghahramani, Zoubin},
  title     = {Dropout as a {B}ayesian Approximation: Representing Model Uncertainty in Deep Learning},
  booktitle = {International Conference on Machine Learning (ICML)},
  year      = {2016}
}

@inproceedings{lee2018,
  author    = {Lee, Kimin and Lee, Kibok and Lee, Honglak and Shin, Jinwoo},
  title     = {A Simple Unified Framework for Detecting Out-of-Distribution Samples and Adversarial Attacks},
  booktitle = {Advances in Neural Information Processing Systems (NeurIPS)},
  year      = {2018}
}

@article{papernot2018,
  author  = {Papernot, Nicolas and McDaniel, Patrick},
  title   = {Deep k-Nearest Neighbors: Towards Confident, Interpretable and Robust Deep Learning},
  journal = {arXiv preprint arXiv:1803.04765},
  year    = {2018}
}

@article{wangmendel1992,
  author  = {Wang, Li-Xin and Mendel, Jerry M.},
  title   = {Generating Fuzzy Rules by Learning from Examples},
  journal = {IEEE Transactions on Systems, Man, and Cybernetics},
  volume  = {22},
  number  = {6},
  pages   = {1414--1427},
  year    = {1992}
}

@article{takagi1985,
  author  = {Takagi, Tomohiro and Sugeno, Michio},
  title   = {Fuzzy Identification of Systems and Its Applications to Modeling and Control},
  journal = {IEEE Transactions on Systems, Man, and Cybernetics},
  volume  = {SMC-15},
  number  = {1},
  pages   = {116--132},
  year    = {1985}
}

@article{tschandl2018,
  author  = {Tschandl, Philipp and Rosendahl, Cliff and Kittler, Harald},
  title   = {The {HAM10000} Dataset, a Large Collection of Multi-Source Dermatoscopic Images of Common Pigmented Skin Lesions},
  journal = {Scientific Data},
  volume  = {5},
  pages   = {180161},
  year    = {2018}
}

@inproceedings{font2025,
  author    = {Font, Miguel N. and Jorro-Aragoneses, Jos{\'e} L. and Ala{\'\i}z, Carlos M.},
  title     = {A Framework for Uncertainty Quantification Based on Nearest Neighbors Across Layers},
  booktitle = {Artificial Neural Networks and Machine Learning -- ICANN 2025},
  series    = {Lecture Notes in Computer Science},
  volume    = {16068},
  pages     = {379--390},
  publisher = {Springer},
  year      = {2026},
  doi       = {10.1007/978-3-032-04558-4_30}
}

@inproceedings{gorbett2022,
  author    = {Gorbett, Matt and Blanchard, Nathaniel},
  title     = {Utilizing Network Features to Detect Erroneous Inputs},
  booktitle = {Proceedings of the IEEE/CVF Winter Conference on Applications of Computer Vision (WACV) Workshops},
  pages     = {34--43},
  year      = {2022}
}

@inproceedings{jaeger2023,
  author    = {Jaeger, Paul F. and L{\"u}th, Carsten T. and Klein, Lukas and Bungert, Till J.},
  title     = {A Call to Reflect on Evaluation Practices for Failure Detection in Image Classification},
  booktitle = {International Conference on Learning Representations (ICLR)},
  year      = {2023}
}

@inproceedings{guo2017,
  author    = {Guo, Chuan and Pleiss, Geoff and Sun, Yu and Weinberger, Kilian Q.},
  title     = {On Calibration of Modern Neural Networks},
  booktitle = {International Conference on Machine Learning (ICML)},
  year      = {2017}
}

@inproceedings{lakshminarayanan2017,
  author    = {Lakshminarayanan, Balaji and Pritzel, Alexander and Blundell, Charles},
  title     = {Simple and Scalable Predictive Uncertainty Estimation Using Deep Ensembles},
  booktitle = {Advances in Neural Information Processing Systems (NeurIPS)},
  year      = {2017}
}

@article{alain2016,
  author  = {Alain, Guillaume and Bengio, Yoshua},
  title   = {Understanding Intermediate Layers Using Linear Classifier Probes},
  journal = {arXiv preprint arXiv:1610.01644},
  year    = {2016}
}

@inproceedings{baldock2021,
  author    = {Baldock, Robert J. N. and Maennel, Hartmut and Neyshabur, Behnam},
  title     = {Deep Learning Through the Lens of Example Difficulty},
  booktitle = {Advances in Neural Information Processing Systems (NeurIPS)},
  year      = {2021}
}

@inproceedings{peterson2019,
  author    = {Peterson, Joshua C. and Battleday, Ruairidh M. and Griffiths, Thomas L. and Russakovsky, Olga},
  title     = {Human Uncertainty Makes Classification More Robust},
  booktitle = {Proceedings of the IEEE/CVF International Conference on Computer Vision (ICCV)},
  year      = {2019}
}

% TODO: conferir volume e páginas no site da IEEE
@article{aghaeipoor2023,
  author  = {Aghaeipoor, Fatemeh and Sabokrou, Mohammad and Fern{\'a}ndez, Alberto},
  title   = {Fuzzy Rule-Based Explainer Systems for Deep Neural Networks: From Local Explainability to Global Understanding},
  journal = {IEEE Transactions on Fuzzy Systems},
  year    = {2023}
}

\end{filecontents*}
% =====================================================================
% Template IJCNN (IEEE conference) — incerteza_fuzzy
% No Overleaf: New Project > Upload Project (zip desta pasta). O IEEEtran.cls já vem no Overleaf.
% Regras usuais da IJCNN (CONFIRA no site da IJCNN 2027 quando a chamada sair):
%   - 6 a 8 páginas COM referências; até 2 extras pagas
%   - revisão duplo-cego: sem nomes, sem agradecimentos, citar trabalho próprio em 3ª pessoa
%   - A4, duas colunas, Times 10 pt, sem número de página
% Números: historico/teste1_2026-10-16 (relatorio.xlsx, RESULTADOS.md), commit f56b043.
% =====================================================================
\documentclass[conference,a4paper]{IEEEtran}
\IEEEoverridecommandlockouts

\usepackage{cite}
\usepackage{amsmath,amssymb,amsfonts}
\usepackage{graphicx}
\usepackage{booktabs,multirow}
\usepackage{textcomp}
\usepackage{xcolor}
\usepackage[hidelinks]{hyperref}

\newcommand{\todo}[1]{\textcolor{red}{[TODO: #1]}}
\newcommand{\term}[1]{\textsc{#1}}

\begin{document}

\title{Reading the Network from the Inside: Interpretable Additive Fuzzy Rules\\
for Detecting Misclassifications of Deep Classifiers}

% ---- versão para submissão (duplo-cego) ----
\author{\IEEEauthorblockN{Anonymous Authors}}
% ---- versão final (camera-ready): descomente e apague a linha acima ----
% \author{\IEEEauthorblockN{Matheus Buttow}
% \IEEEauthorblockA{\textit{Programa de Pós-Graduação em ...} \\
% \textit{Universidade ...}\\
% Cidade, Brasil \\
% matheusbuttow@gmail.com}}

\maketitle

\begin{abstract}
Deep classifiers often fail with high confidence, which is critical in medical applications such as skin-cancer screening,
where an undetected error delays a diagnosis.
We propose an interpretable detector of misclassifications that reads the internal behavior of a trained, frozen network:
linear probes at every block reveal whether the network changed its mind on the way to the answer,
augmented views reveal whether the answer is stable, and nearest neighbors in feature space reveal whether the input is familiar.
These signals are combined by an additive fuzzy model, a zero-order Takagi--Sugeno system with single-antecedent rules and logistic output,
so that every rule has a readable form (``IF final-block agreement is \term{low} THEN the risk of error increases'') and a learned weight.
On HAM10000 (benign vs.\ malignant, EfficientNet-B0) and six additional scenarios (CIFAR-10 with and without 30\% label noise,
CIFAR-100, SVHN and MNIST), all evaluated on the same frozen networks and test images as twelve competing detectors,
the method obtains the best or tied-best mean AUROC in seven of eight settings and stays within 0.004 of the best in the remaining one
(AUROC 0.916 on HAM10000 and 0.903 on CIFAR-10 with label noise, against 0.865 and 0.797 for the softmax confidence).
Against the probe-based meta-model of Chen et al., it is significantly better in 18 of 26 runs and never significantly worse,
while its rules remain stable across resamplings.
\end{abstract}

\begin{IEEEkeywords}
uncertainty estimation, misclassification detection, fuzzy systems, interpretability, skin lesion classification
\end{IEEEkeywords}

% =====================================================================
\section{Introduction}

Deep networks for dermoscopy already reach expert-level accuracy on curated benchmarks, yet they still make mistakes, and some of these
mistakes come with high confidence. In a screening workflow, the cost is asymmetric: a malignant lesion classified as benign with a
confident score is unlikely to be reviewed. A practical way to deploy such a model is \emph{selective prediction}: the network answers
the cases it is likely to get right and refers the rest to a dermatologist. This requires a score that ranks images by their risk of
being misclassified, a task known as misclassification (or failure) detection~\cite{hendrycks2017,jaeger2023}.

The simplest score is the network's own confidence, the maximum softmax probability~\cite{hendrycks2017}. It is a strong baseline,
but it is poorly calibrated~\cite{guo2017} and, by construction, it sees only the last layer. Bayesian approximations such as
MC Dropout~\cite{gal2016} and test-time augmentation still summarize the output, and distance-based scores such as
Mahalanobis~\cite{lee2018} and Deep $k$NN~\cite{papernot2018} target out-of-distribution inputs rather than in-distribution errors.
Meta-models that read the internal representations, such as the linear-probe approach of Chen et al.~\cite{chen2019},
layer-wise neighbors~\cite{font2025} or classifiers on penultimate features~\cite{gorbett2022}, are more informative,
but they are black boxes themselves: they say \emph{that} an image is risky, not \emph{why}.

We argue that the reasons a network fails are visible in how it reaches its answer, and that they can be stated as rules.
A correct, confident prediction is usually formed early and kept until the end; a wrong one tends to hesitate across blocks,
to change when the image is mirrored, or to sit among training neighbors that received a different answer.
We measure these behaviors with a small set of features and combine them with an additive fuzzy model whose rules a clinician can read.

Our contributions are:
\begin{itemize}
  \item a reading of a frozen network, without retraining, along four complementary axes (the answer, the path through the blocks,
        the stability across views and the familiarity of the input), summarized in a few interpretable features;
  \item an additive fuzzy detector, a zero-order Takagi--Sugeno system with single-antecedent rules and logistic output, whose rule
        weights are stable across resamplings (Spearman correlation 0.91--0.94 on HAM10000);
  \item a controlled comparison with twelve detectors on eight scenarios, using the same frozen networks, the same test images and the
        same cross-validation folds, with paired bootstrap tests, plus three checks against real-world uncertainty
        (biopsy-confirmed lesions, multiple photographs of the same lesion and human disagreement on CIFAR-10H~\cite{peterson2019}).
\end{itemize}

% =====================================================================
\section{Related Work}

\subsection{Misclassification and failure detection}
The maximum softmax probability (MSP)~\cite{hendrycks2017} remains the reference baseline, and a large benchmark showed that most
dedicated methods fail to beat it consistently on in-distribution errors~\cite{jaeger2023}. MC Dropout~\cite{gal2016} and deep
ensembles~\cite{lakshminarayanan2017} estimate epistemic uncertainty at a higher inference cost. ConfidNet~\cite{corbiere2019}
learns to regress the true-class probability from penultimate features. Distance-based scores, Mahalanobis~\cite{lee2018} and Deep
$k$NN~\cite{papernot2018}, measure familiarity in feature space and are mostly designed for out-of-distribution detection.

\subsection{Meta-models on internal representations}
Linear probes~\cite{alain2016} show how linearly decodable the label is at each layer. Chen et al.~\cite{chen2019} train one probe per
layer and feed all probe outputs to a gradient-boosting meta-model that predicts whether the base network is correct; their gains are
largest under label noise. Font et al.~\cite{font2025} compare the nearest training neighbors across layers, and Gorbett and
Blanchard~\cite{gorbett2022} train a classifier on penultimate features. The depth at which a prediction is formed has also been related
to example difficulty~\cite{baldock2021}. These works confirm that the internal signal exists, but their meta-models are not interpretable.

\subsection{Fuzzy systems for uncertainty and explanation}
Fuzzy rule-based systems have been used to explain deep networks from attribution values~\cite{aghaeipoor2023} and to model uncertainty
inside the network through type-2 or neuro-fuzzy layers. Rule learning from data dates back to Wang and Mendel~\cite{wangmendel1992},
and Takagi--Sugeno models~\cite{takagi1985} combine linguistic antecedents with numeric consequents. To our knowledge, no previous work
applies fuzzy rules to the internal behavior of a trained network to detect its errors.

% =====================================================================
\section{Method}

Let $f$ be a trained classifier with $K$ classes, kept frozen. For an input $x$, $f$ predicts $\hat y(x)$ with confidence $p(x)$:
in the binary case (HAM10000), $p$ is the sigmoid output $P(\text{malignant}\mid x)$ and $\hat y = \mathbf{1}[p \ge \tau]$, with the
threshold $\tau$ chosen on a validation set (Youden index); in the multiclass case, $p$ is the maximum softmax probability.
The goal is a score $r(x)$ that is high when $\hat y(x)$ is wrong.

\subsection{Reading the network}
\label{sec:reading}

\textit{The path through the blocks.}
Let $a_b(x)$ be the output of block $b \in \{1,\dots,B\}$, averaged over the spatial positions. On a set of reference training images,
we train one probe per block, $g_b$, a logistic regression (with standardization, and PCA to 256 dimensions for wide layers) whose target
is the \emph{network's decision} $\hat y$, not the label. The probe therefore answers: ``had the network stopped at block $b$, what would
it answer?''. For a test image, let $c_b(x) = g_b(\hat y(x) \mid a_b(x))$ be the probability that block $b$ already gives the final answer,
and $s_b(x) = \mathbf{1}[\arg\max g_b(a_b(x)) = \hat y(x)]$ its agreement. We compute
\begin{align}
  \textit{switches} &= \textstyle\sum_{b=2}^{B} |s_b - s_{b-1}|, \\
  \textit{decision depth} &= \tfrac{1}{B}\max\{b : s_b = 0\}, \\
  \textit{probe entropy} &= \tfrac{1}{B}\textstyle\sum_{b} H\big(g_b(a_b)\big)/\log K, \\
  \textit{final agreement} &= \tfrac{1}{|L|}\textstyle\sum_{b \in L} c_b, \\
  \textit{final jump} &= |p(x) - c_B(x)|.
\end{align}
where $L$ holds the last $\lfloor B/3 \rfloor$ blocks; the decision depth is zero when all blocks agree with $\hat y$ and follows the prediction depth of Baldock et al.~\cite{baldock2021}.

\textit{Stability across views.}
Each image is also passed through $V-1$ label-preserving transformations (flips and a 180$^\circ$ rotation for skin lesions; a horizontal flip and
2-pixel shifts for CIFAR; small shifts only for SVHN and MNIST, where flips change the digit). We record the fraction of views whose decision differs from $\hat y$ (\textit{view flips}),
the standard deviation of the confidence in $\hat y$ across views (\textit{view std}) and the standard deviation of the probe path
$(c_1,\dots,c_B)$ across views, averaged over blocks (\textit{path std}).

\textit{Familiarity.}
In the penultimate space, we retrieve the $k=15$ nearest reference images (cosine similarity) and record the fraction whose network
decision differs from $\hat y$ (\textit{neighbor disagreement}) and their mean distance divided by the median distance among reference
images (\textit{neighbor distance}).

Together with $p$ itself, these eleven features describe the answer, the path, the stability and the familiarity of each prediction.
The three view features require $V$ forward passes; all others require one pass.

\subsection{Additive fuzzy model}
Each selected feature $x_j$ is covered by three triangular membership functions, \term{low}, \term{medium} and \term{high}, forming a strong
partition with peaks at its 5th, 50th and 95th percentiles. For the binary output $p$, the peaks are fixed at $\{0, \tau, 1\}$, so that
\term{medium} means ``at the decision threshold''. The risk of error is
\begin{equation}
  r(\mathbf{x}) = \sigma\Big(b + \sum_{j=1}^{d}\sum_{k=1}^{3} w_{jk}\,\mu_{jk}(x_j)\Big),
  \label{eq:aditivo}
\end{equation}
where $\mu_{jk}$ is the membership of $x_j$ in term $k$ and $\sigma$ is the logistic function. Each pair $(j,k)$ is a rule
``IF $x_j$ is $k$ THEN the log-odds of error change by $w_{jk}$'', so the model is a zero-order Takagi--Sugeno system~\cite{takagi1985}
with single-antecedent rules whose consequents are added and passed through $\sigma$. The weights are learned by $\ell_2$-regularized
logistic regression on the memberships, with the network's error ($\hat y \ne y$) as target. With $d=4$ features the model has 12 rules
and a bias; a positive weight means that the condition raises the risk.

\subsection{Feature selection}
We always include $p$, since the answer itself is the strongest single signal, and add the $d-1=3$ features with the highest
\emph{AUROC beyond $P$}: the AUROC for detecting errors computed within five quantile bins of $p$ and averaged with bin sizes as weights.
It measures what a feature adds among images with similar confidence. A candidate is skipped if its absolute Spearman correlation with
an already selected feature is at least 0.85. Selection is repeated inside every training fold. A variant without the three view features
(one forward pass) is also reported.

% =====================================================================
\section{Experimental Setup}

\subsection{Scenarios and networks}
Table~\ref{tab:redes} lists the eight scenarios. HAM10000~\cite{tschandl2018} is split by lesion (one image per lesion, so no lesion
appears in two sets) into training, validation and test sets (70/10/20\%), with \textit{mel}, \textit{bcc} and \textit{akiec} as malignant.
We use all lesions (five seeds) and a 3,000-lesion subset (three seeds). An ImageNet-pretrained EfficientNet-B0 with a two-layer head is
fine-tuned for 15 epochs (AdamW), and the epoch and threshold are selected on the validation set; its 18 sub-blocks are probed.
The other scenarios reproduce the settings of the papers we compare with: ResNet-32 on CIFAR-10, clean and with 30\% of the training
labels replaced by another class~\cite{chen2019}; VGG-16 on CIFAR-10 and CIFAR-100, a small CNN on SVHN and an MLP on MNIST~\cite{corbiere2019,jaeger2023}.
Each network is trained once per seed and kept frozen; 5,000 training images serve as reference for probes, neighbors and Mahalanobis.

\subsection{Competing detectors}
All detectors score the same test images of the same frozen network (higher score = more likely error).
\emph{Without training:} softmax confidence (MSP) and entropy~\cite{hendrycks2017}, distance to the threshold (binary), view standard
deviation (TTA), MC Dropout with $T=30$ passes through the head~\cite{gal2016}, Mahalanobis with a shared covariance~\cite{lee2018}
and Deep $k$NN~\cite{papernot2018}.
\emph{Trained:} an output-only model (logistic regression on $p$ and $p^2$ in the binary case; equal to MSP in the multiclass case),
Chen et al.~\cite{chen2019} (label-trained probes at every block, top-3 probabilities per probe, gradient boosting),
Font et al.~\cite{font2025} (layer-wise nearest neighbors), Gorbett and Blanchard~\cite{gorbett2022} (linear SVM on penultimate features)
and a post-hoc version of ConfidNet~\cite{corbiere2019} (an MLP that regresses the true-class probability from penultimate features).
All trained detectors, including ours, use the same stratified 5-fold cross-validation on the test set, and are evaluated on their
out-of-fold scores.

\subsection{Metrics}
The main metric is the AUROC for separating the network's errors from its correct predictions. We also report the AUPR of errors,
FPR at 95\% TPR, the excess area under the risk--coverage curve (E-AURC)~\cite{jaeger2023} and the number of errors found when the
10\%, 20\% and 30\% most uncertain images are referred to an expert. Differences in AUROC are tested per run with a paired bootstrap
(1,000 resamples shared by all methods); a difference is significant when its 95\% confidence interval excludes zero.

% =====================================================================
\section{Results}

% Gerado por tabelas_latex.py a partir de relatorio.xlsx.
% Preâmbulo: \usepackage{booktabs,multirow,graphicx}

% ==================================================================
% Tabela 1: Redes
% ==================================================================
\begin{table}[t]
\centering
\caption{Evaluated (frozen) networks. Mean over seeds.}
\label{tab:redes}
\resizebox{\columnwidth}{!}{%
\begin{tabular}{llcccc}
  \toprule
  Scenario & Network & Classes & \begin{tabular}[c]{@{}c@{}}Test\\images\end{tabular} & Accuracy & Seeds \\
  \midrule
  HAM10000 (all) & EfficientNet-B0 & 2 & 1,495 & 0.861 & 5 \\
  HAM10000 (3000) & EfficientNet-B0 & 2 & 601 & 0.831 & 3 \\
  CIFAR-10 30\% noise & ResNet-32 & 10 & 10,000 & 0.865 & 3 \\
  CIFAR-100 & VGG-16 & 100 & 10,000 & 0.731 & 3 \\
  CIFAR-10 (ResNet) & ResNet-32 & 10 & 10,000 & 0.929 & 3 \\
  CIFAR-10 (VGG) & VGG-16 & 10 & 10,000 & 0.937 & 3 \\
  SVHN & CNN & 10 & 26,032 & 0.954 & 3 \\
  MNIST & MLP & 10 & 10,000 & 0.991 & 3 \\
  \bottomrule
\end{tabular}}
\end{table}

% ==================================================================
% Tabela 2: AUROC
% ==================================================================
\begin{table*}[t]
\centering
\caption{AUROC for detecting the network's errors (mean $\pm$ std over seeds). Best mean per scenario in bold.}
\label{tab:auroc}
\resizebox{\textwidth}{!}{%
\begin{tabular}{lcccccccc}
  \toprule
  Method & \begin{tabular}[c]{@{}c@{}}HAM10000\\(all)\end{tabular} & \begin{tabular}[c]{@{}c@{}}HAM10000\\(3000)\end{tabular} & \begin{tabular}[c]{@{}c@{}}CIFAR-10 30\%\\noise\end{tabular} & CIFAR-100 & \begin{tabular}[c]{@{}c@{}}CIFAR-10\\(ResNet)\end{tabular} & \begin{tabular}[c]{@{}c@{}}CIFAR-10\\(VGG)\end{tabular} & SVHN & MNIST \\
  \midrule
  Softmax (MSP) & 0.865\,{\scriptsize$\pm$0.038} & 0.819\,{\scriptsize$\pm$0.099} & 0.797\,{\scriptsize$\pm$0.014} & 0.866\,{\scriptsize$\pm$0.003} & 0.928\,{\scriptsize$\pm$0.004} & 0.915\,{\scriptsize$\pm$0.004} & \textbf{0.926}\,{\scriptsize$\pm$0.001} & 0.985\,{\scriptsize$\pm$0.000} \\
  Softmax entropy & 0.865\,{\scriptsize$\pm$0.038} & 0.819\,{\scriptsize$\pm$0.099} & 0.778\,{\scriptsize$\pm$0.015} & 0.867\,{\scriptsize$\pm$0.003} & 0.928\,{\scriptsize$\pm$0.004} & 0.914\,{\scriptsize$\pm$0.004} & 0.924\,{\scriptsize$\pm$0.001} & 0.984\,{\scriptsize$\pm$0.000} \\
  Distance to threshold & 0.862\,{\scriptsize$\pm$0.013} & 0.842\,{\scriptsize$\pm$0.069} & -- & -- & -- & -- & -- & -- \\
  TTA & 0.864\,{\scriptsize$\pm$0.019} & 0.804\,{\scriptsize$\pm$0.067} & 0.810\,{\scriptsize$\pm$0.012} & 0.845\,{\scriptsize$\pm$0.002} & 0.931\,{\scriptsize$\pm$0.003} & 0.931\,{\scriptsize$\pm$0.005} & 0.909\,{\scriptsize$\pm$0.003} & \textbf{0.986}\,{\scriptsize$\pm$0.000} \\
  MC Dropout & 0.864\,{\scriptsize$\pm$0.039} & 0.818\,{\scriptsize$\pm$0.100} & 0.779\,{\scriptsize$\pm$0.013} & 0.864\,{\scriptsize$\pm$0.002} & 0.926\,{\scriptsize$\pm$0.004} & 0.910\,{\scriptsize$\pm$0.007} & 0.918\,{\scriptsize$\pm$0.002} & 0.984\,{\scriptsize$\pm$0.001} \\
  Mahalanobis & 0.178\,{\scriptsize$\pm$0.028} & 0.214\,{\scriptsize$\pm$0.091} & 0.695\,{\scriptsize$\pm$0.014} & 0.847\,{\scriptsize$\pm$0.002} & 0.802\,{\scriptsize$\pm$0.009} & 0.918\,{\scriptsize$\pm$0.002} & 0.492\,{\scriptsize$\pm$0.009} & 0.769\,{\scriptsize$\pm$0.015} \\
  Deep kNN & 0.610\,{\scriptsize$\pm$0.012} & 0.585\,{\scriptsize$\pm$0.070} & 0.699\,{\scriptsize$\pm$0.008} & 0.667\,{\scriptsize$\pm$0.002} & 0.738\,{\scriptsize$\pm$0.004} & 0.594\,{\scriptsize$\pm$0.013} & 0.859\,{\scriptsize$\pm$0.005} & 0.888\,{\scriptsize$\pm$0.007} \\
  \midrule
  Output only$^\dagger$ & 0.907\,{\scriptsize$\pm$0.017} & 0.872\,{\scriptsize$\pm$0.060} & 0.797\,{\scriptsize$\pm$0.014} & 0.866\,{\scriptsize$\pm$0.003} & 0.928\,{\scriptsize$\pm$0.004} & 0.915\,{\scriptsize$\pm$0.004} & \textbf{0.926}\,{\scriptsize$\pm$0.001} & 0.985\,{\scriptsize$\pm$0.000} \\
  Chen et al.~\cite{chen2019} & 0.909\,{\scriptsize$\pm$0.017} & 0.896\,{\scriptsize$\pm$0.031} & 0.863\,{\scriptsize$\pm$0.005} & 0.871\,{\scriptsize$\pm$0.003} & 0.920\,{\scriptsize$\pm$0.004} & 0.918\,{\scriptsize$\pm$0.002} & 0.921\,{\scriptsize$\pm$0.004} & 0.976\,{\scriptsize$\pm$0.004} \\
  Font et al.~\cite{font2025} & 0.907\,{\scriptsize$\pm$0.019} & 0.886\,{\scriptsize$\pm$0.044} & 0.820\,{\scriptsize$\pm$0.011} & 0.847\,{\scriptsize$\pm$0.003} & 0.886\,{\scriptsize$\pm$0.007} & 0.865\,{\scriptsize$\pm$0.005} & 0.919\,{\scriptsize$\pm$0.003} & 0.976\,{\scriptsize$\pm$0.001} \\
  Gorbett and Blanchard~\cite{gorbett2022} & 0.841\,{\scriptsize$\pm$0.042} & 0.837\,{\scriptsize$\pm$0.078} & 0.880\,{\scriptsize$\pm$0.005} & 0.858\,{\scriptsize$\pm$0.003} & 0.921\,{\scriptsize$\pm$0.006} & 0.922\,{\scriptsize$\pm$0.003} & 0.913\,{\scriptsize$\pm$0.002} & 0.932\,{\scriptsize$\pm$0.031} \\
  ConfidNet~\cite{corbiere2019} & 0.790\,{\scriptsize$\pm$0.032} & 0.757\,{\scriptsize$\pm$0.078} & 0.841\,{\scriptsize$\pm$0.007} & 0.816\,{\scriptsize$\pm$0.003} & 0.901\,{\scriptsize$\pm$0.007} & 0.876\,{\scriptsize$\pm$0.012} & 0.853\,{\scriptsize$\pm$0.006} & 0.792\,{\scriptsize$\pm$0.028} \\
  \midrule
  Additive fuzzy, no views (ours) & 0.914\,{\scriptsize$\pm$0.016} & 0.901\,{\scriptsize$\pm$0.048} & 0.896\,{\scriptsize$\pm$0.004} & 0.875\,{\scriptsize$\pm$0.002} & 0.919\,{\scriptsize$\pm$0.004} & 0.922\,{\scriptsize$\pm$0.004} & 0.925\,{\scriptsize$\pm$0.000} & 0.981\,{\scriptsize$\pm$0.000} \\
  \textbf{Additive fuzzy (ours)} & \textbf{0.916}\,{\scriptsize$\pm$0.012} & \textbf{0.904}\,{\scriptsize$\pm$0.044} & \textbf{0.903}\,{\scriptsize$\pm$0.003} & \textbf{0.886}\,{\scriptsize$\pm$0.002} & \textbf{0.936}\,{\scriptsize$\pm$0.002} & \textbf{0.934}\,{\scriptsize$\pm$0.002} & \textbf{0.926}\,{\scriptsize$\pm$0.002} & 0.982\,{\scriptsize$\pm$0.000} \\
  \bottomrule
\end{tabular}}

\vspace{2pt}
{\footnotesize $^\dagger$ Binary: logistic regression on $P$ and $P^2$; multiclass: identical to softmax.}
\end{table*}

% ==================================================================
% Tabela 3: Erros encontrados
% ==================================================================
\begin{table*}[t]
\centering
\caption{Network errors found when the 10\%, 20\% and 30\% most uncertain images are referred to an expert (per 1,000 images; MNIST per 10,000). In parentheses, the fraction of all network errors.}
\label{tab:encontrados}
\begin{minipage}[t]{0.495\textwidth}
\centering
\scriptsize\setlength{\tabcolsep}{3pt}
\begin{tabular}[t]{lccc}
  \toprule
  Method & \begin{tabular}[c]{@{}c@{}}Reviewing\\10\%\end{tabular} & \begin{tabular}[c]{@{}c@{}}Reviewing\\20\%\end{tabular} & \begin{tabular}[c]{@{}c@{}}Reviewing\\30\%\end{tabular} \\
  \midrule
  \multicolumn{4}{@{}l}{\textit{HAM10000 (all)} --- 139 errors} \\
  Random & 14 (10\%) & 28 (20\%) & 42 (30\%) \\
  Softmax (MSP) & 49 (35\%) & 82 (59\%) & 111 (80\%) \\
  MC Dropout & 48 (34\%) & 82 (59\%) & 111 (80\%) \\
  TTA & 49 (35\%) & 87 (62\%) & 111 (80\%) \\
  Chen et al.~\cite{chen2019} & 61 (44\%) & 104 (75\%) & 122 (88\%) \\
  Gorbett and Blanchard~\cite{gorbett2022} & 54 (39\%) & 88 (64\%) & 110 (79\%) \\
  Output only & \textbf{65 (47\%)} & \textbf{107 (77\%)} & 125 (90\%) \\
  Additive fuzzy, no views (ours) & 63 (45\%) & 106 (76\%) & 125 (90\%) \\
  \textbf{Additive fuzzy (ours)} & 64 (46\%) & 106 (76\%) & \textbf{126 (91\%)} \\
  \midrule
  \multicolumn{4}{@{}l}{\textit{HAM10000 (3000)} --- 169 errors} \\
  Random & 17 (10\%) & 34 (20\%) & 51 (30\%) \\
  Softmax (MSP) & 34 (20\%) & 82 (48\%) & 126 (74\%) \\
  MC Dropout & 34 (20\%) & 80 (48\%) & 126 (74\%) \\
  TTA & 37 (22\%) & 74 (44\%) & 113 (67\%) \\
  Chen et al.~\cite{chen2019} & 58 (34\%) & 109 (64\%) & 147 (87\%) \\
  Gorbett and Blanchard~\cite{gorbett2022} & 55 (33\%) & 101 (60\%) & 131 (77\%) \\
  Output only & 44 (26\%) & 109 (64\%) & \textbf{148 (87\%)} \\
  Additive fuzzy, no views (ours) & 57 (33\%) & \textbf{119 (70\%)} & \textbf{148 (88\%)} \\
  \textbf{Additive fuzzy (ours)} & \textbf{63 (37\%)} & 118 (70\%) & \textbf{148 (87\%)} \\
  \midrule
  \multicolumn{4}{@{}l}{\textit{CIFAR-10 30\% noise} --- 135 errors} \\
  Random & 13 (10\%) & 27 (20\%) & 40 (30\%) \\
  Softmax (MSP) = Output only & 53 (40\%) & 84 (62\%) & 98 (73\%) \\
  MC Dropout & 49 (36\%) & 78 (58\%) & 94 (70\%) \\
  TTA & 38 (28\%) & 71 (53\%) & 96 (71\%) \\
  Chen et al.~\cite{chen2019} & 53 (39\%) & 88 (66\%) & 108 (80\%) \\
  Gorbett and Blanchard~\cite{gorbett2022} & 56 (41\%) & 93 (69\%) & 113 (84\%) \\
  Additive fuzzy, no views (ours) & \textbf{62 (46\%)} & 96 (72\%) & 116 (86\%) \\
  \textbf{Additive fuzzy (ours)} & \textbf{62 (46\%)} & \textbf{99 (74\%)} & \textbf{119 (88\%)} \\
  \midrule
  \multicolumn{4}{@{}l}{\textit{CIFAR-100} --- 269 errors} \\
  Random & 27 (10\%) & 54 (20\%) & 81 (30\%) \\
  Softmax (MSP) = Output only & 74 (28\%) & 136 (50\%) & 187 (70\%) \\
  MC Dropout & 77 (29\%) & 138 (51\%) & 186 (69\%) \\
  TTA & 53 (20\%) & 114 (42\%) & 176 (65\%) \\
  Chen et al.~\cite{chen2019} & 76 (28\%) & 137 (51\%) & 187 (69\%) \\
  Gorbett and Blanchard~\cite{gorbett2022} & 74 (28\%) & 136 (50\%) & 184 (68\%) \\
  Additive fuzzy, no views (ours) & 78 (29\%) & 140 (52\%) & 189 (70\%) \\
  \textbf{Additive fuzzy (ours)} & \textbf{81 (30\%)} & \textbf{145 (54\%)} & \textbf{194 (72\%)} \\
  \bottomrule
\end{tabular}
\end{minipage}\hfill
\begin{minipage}[t]{0.495\textwidth}
\centering
\scriptsize\setlength{\tabcolsep}{3pt}
\begin{tabular}[t]{lccc}
  \toprule
  Method & \begin{tabular}[c]{@{}c@{}}Reviewing\\10\%\end{tabular} & \begin{tabular}[c]{@{}c@{}}Reviewing\\20\%\end{tabular} & \begin{tabular}[c]{@{}c@{}}Reviewing\\30\%\end{tabular} \\
  \midrule
  \multicolumn{4}{@{}l}{\textit{CIFAR-10 (ResNet)} --- 71 errors} \\
  Random & 7 (10\%) & 14 (20\%) & 21 (30\%) \\
  Softmax (MSP) = Output only & 44 (62\%) & 61 (87\%) & 67 (95\%) \\
  MC Dropout & 44 (61\%) & 61 (87\%) & 67 (95\%) \\
  TTA & 45 (63\%) & \textbf{64 (90\%)} & \textbf{68 (96\%)} \\
  Chen et al.~\cite{chen2019} & 42 (60\%) & 60 (84\%) & 67 (94\%) \\
  Gorbett and Blanchard~\cite{gorbett2022} & 43 (61\%) & 60 (85\%) & 66 (93\%) \\
  Additive fuzzy, no views (ours) & 43 (61\%) & 59 (84\%) & 66 (93\%) \\
  \textbf{Additive fuzzy (ours)} & \textbf{48 (68\%)} & \textbf{64 (90\%)} & 67 (95\%) \\
  \midrule
  \multicolumn{4}{@{}l}{\textit{CIFAR-10 (VGG)} --- 63 errors} \\
  Random & 6 (10\%) & 13 (20\%) & 19 (30\%) \\
  Softmax (MSP) = Output only & 41 (65\%) & 54 (87\%) & 58 (93\%) \\
  MC Dropout & 40 (64\%) & 54 (86\%) & 57 (92\%) \\
  TTA & 43 (69\%) & \textbf{57 (90\%)} & \textbf{60 (96\%)} \\
  Chen et al.~\cite{chen2019} & 37 (59\%) & 54 (86\%) & 59 (94\%) \\
  Gorbett and Blanchard~\cite{gorbett2022} & 39 (62\%) & 54 (86\%) & \textbf{60 (95\%)} \\
  Additive fuzzy, no views (ours) & 40 (65\%) & 54 (86\%) & 59 (94\%) \\
  \textbf{Additive fuzzy (ours)} & \textbf{44 (70\%)} & 56 (90\%) & 59 (95\%) \\
  \midrule
  \multicolumn{4}{@{}l}{\textit{SVHN} --- 46 errors} \\
  Random & 5 (10\%) & 9 (20\%) & 14 (30\%) \\
  Softmax (MSP) = Output only & \textbf{35 (76\%)} & 41 (90\%) & \textbf{43 (94\%)} \\
  MC Dropout & 32 (69\%) & 41 (89\%) & \textbf{43 (93\%)} \\
  TTA & 30 (66\%) & 41 (89\%) & \textbf{43 (94\%)} \\
  Chen et al.~\cite{chen2019} & 34 (75\%) & 41 (89\%) & \textbf{43 (94\%)} \\
  Gorbett and Blanchard~\cite{gorbett2022} & \textbf{35 (76\%)} & 41 (88\%) & 42 (91\%) \\
  Additive fuzzy, no views (ours) & \textbf{35 (77\%)} & 41 (90\%) & \textbf{43 (94\%)} \\
  \textbf{Additive fuzzy (ours)} & \textbf{35 (76\%)} & \textbf{42 (90\%)} & \textbf{43 (94\%)} \\
  \midrule
  \multicolumn{4}{@{}l}{\textit{MNIST} --- 93 errors (per 10,000)} \\
  Random & 9 (10\%) & 19 (20\%) & 28 (30\%) \\
  Softmax (MSP) = Output only & 92 (99\%) & \textbf{93 (100\%)} & \textbf{93 (100\%)} \\
  MC Dropout & 92 (99\%) & \textbf{93 (100\%)} & \textbf{93 (100\%)} \\
  TTA & \textbf{93 (100\%)} & \textbf{93 (100\%)} & \textbf{93 (100\%)} \\
  Chen et al.~\cite{chen2019} & 91 (97\%) & \textbf{93 (100\%)} & \textbf{93 (100\%)} \\
  Gorbett and Blanchard~\cite{gorbett2022} & 79 (85\%) & 84 (90\%) & 87 (94\%) \\
  Additive fuzzy, no views (ours) & 92 (99\%) & \textbf{93 (100\%)} & \textbf{93 (100\%)} \\
  \textbf{Additive fuzzy (ours)} & \textbf{93 (100\%)} & \textbf{93 (100\%)} & \textbf{93 (100\%)} \\
  \bottomrule
\end{tabular}
\end{minipage}
\end{table*}

% ==================================================================
% Tabela 4: Ganho sobre o softmax
% ==================================================================
\begin{table}[t]
\centering
\caption{AUROC gain of the additive fuzzy model over softmax per scenario.}
\label{tab:ganho}
\resizebox{\columnwidth}{!}{%
\begin{tabular}{lccccc}
  \toprule
  Scenario & Classes & \begin{tabular}[c]{@{}c@{}}Error\\rate\end{tabular} & Softmax & \begin{tabular}[c]{@{}c@{}}Additive\\fuzzy\end{tabular} & Gain \\
  \midrule
  MNIST & 10 & 0.9\% & 0.985 & 0.982 & $-$0.003 \\
  SVHN & 10 & 4.6\% & 0.926 & 0.926 & $-$0.001 \\
  CIFAR-10 (ResNet) & 10 & 7.1\% & 0.928 & 0.936 & +0.008 \\
  CIFAR-10 (VGG) & 10 & 6.3\% & 0.915 & 0.934 & +0.019 \\
  CIFAR-100 & 100 & 26.9\% & 0.866 & 0.886 & +0.020 \\
  HAM10000 (all) & 2 & 13.9\% & 0.865 & 0.916 & +0.052 \\
  HAM10000 (3000) & 2 & 16.9\% & 0.819 & 0.904 & +0.085 \\
  CIFAR-10 30\% noise & 10 & 13.5\% & 0.797 & 0.903 & +0.106 \\
  \bottomrule
\end{tabular}}
\end{table}

% ==================================================================
% Tabela 5: Significância
% ==================================================================
\begin{table}[t]
\centering
\caption{Runs (seeds) in which the additive fuzzy model is significantly better (95\% CI, paired bootstrap) than Chen et al. and the output-only baseline.}
\label{tab:signif}
\resizebox{\columnwidth}{!}{%
\begin{tabular}{lcccc}
  \toprule
  Scenario & \begin{tabular}[c]{@{}c@{}}Beats\\Chen\end{tabular} & \begin{tabular}[c]{@{}c@{}}Loses to\\Chen\end{tabular} & \begin{tabular}[c]{@{}c@{}}Beats output\\only\end{tabular} & Seeds \\
  \midrule
  HAM10000 (all) & 1 & 0 & 2 & 5 \\
  HAM10000 (3000) & 1 & 0 & 2 & 3 \\
  CIFAR-10 30\% noise & 3 & 0 & 3 & 3 \\
  CIFAR-100 & 3 & 0 & 3 & 3 \\
  CIFAR-10 (ResNet) & 3 & 0 & 3 & 3 \\
  CIFAR-10 (VGG) & 3 & 0 & 3 & 3 \\
  SVHN & 2 & 0 & 0 & 3 \\
  MNIST & 2 & 0 & 0 & 3 \\
  \midrule
  \textbf{Total} & \textbf{18} & \textbf{0} & \textbf{16} & \textbf{26} \\
  \bottomrule
\end{tabular}}
\end{table}


\subsection{Detection performance}
Table~\ref{tab:auroc} reports the AUROC. The additive fuzzy model has the best mean in six of the eight scenarios and ties the softmax
on SVHN (0.926); on MNIST, where the network errs in under 1\% of the images, it is 0.004 below the best detector. It also has the lowest
E-AURC in the same six scenarios, e.g.\ 14.8 vs.\ 22.0 for Chen et al.\ and 45.4 for MSP on noisy CIFAR-10.
Table~\ref{tab:ganho} shows that the gain over the softmax grows when the softmax is less reliable: it is negligible on MNIST and SVHN,
0.01--0.02 on clean CIFAR, 0.05--0.09 on HAM10000 and 0.11 on CIFAR-10 with label noise. Against Chen et al., the closest learned competitor,
the fuzzy model is significantly better in 18 of 26 runs and significantly worse in none (Table~\ref{tab:signif}).

On HAM10000 the picture is more balanced. The fuzzy model (0.916), Chen et al.\ (0.909) and the output-only model (0.907) are close,
and the gain over Chen et al.\ is significant in only one of five seeds. When 20\% of the lesions are referred, the three methods find
106, 104 and 107 of every 139 errors per 1,000 lesions (Table~\ref{tab:encontrados}), against 82 for MSP and MC Dropout. In this
scenario, the advantage of the fuzzy model is its explanation, not its ranking. Mahalanobis is inverted on HAM10000 (AUROC 0.18):
most errors are typical-looking lesions, close to the training data, rather than unusual ones.

The single-pass variant, without views, keeps most of the gain where it matters (0.914 on HAM10000, 0.896 on noisy CIFAR-10), but on
clean CIFAR-10 the views are needed to beat the softmax (0.919 vs.\ 0.928).

\subsection{Rules and stability}
The rules are short and consistent with intuition. On HAM10000 (seed 42), being at the threshold raises the log-odds of error
(\term{medium} $p$: $+1.32$), a low final jump lowers it ($-1.53$) and a high path std across views raises it ($+1.16$).
On noisy CIFAR-10, the dominant rules concern the end of the path: \term{low} final agreement $+1.57$, \term{high} final agreement $-2.18$.
Across ten bootstrap resamples, the Spearman correlation between rule weights is 0.91--0.94 on HAM10000 (all lesions) and 0.90--0.98
on the CIFAR scenarios. Fig.~\ref{fig:pertinencias} shows the membership functions and weights.

\begin{figure}[t]
\centering
\IfFileExists{figuras/pertinencias.png}{\includegraphics[width=\columnwidth]{figuras/pertinencias.png}}{\fbox{\parbox{0.9\columnwidth}{\centering\vspace{1cm}
coloque figuras/pertinencias.png (gerada por mostrar\_pertinencias.py)\vspace{1cm}}}}
\caption{Membership functions (top) and rule weights (bottom) of the additive fuzzy model on HAM10000. A positive weight raises the risk of error.}
\label{fig:pertinencias}
\end{figure}

\subsection{Checks against real-world uncertainty}
\textit{Biopsy.} On HAM10000 (all lesions, five seeds), the network errs on 29.1\% of the histopathology-confirmed lesions and on 3.1\%
of those confirmed by follow-up, consensus or confocal microscopy: the cases on which the network hesitates are the ones that clinicians
also sent to biopsy.
\textit{Other photographs of the same lesion.} HAM10000 contains extra photographs of many test lesions, never used in training. The network
changes its answer on 21\% of these lesions. Here, the distance to the threshold predicts this change best (AUROC 0.72), and the fuzzy model
is less effective (0.63): it is tuned to predict errors, not instability across acquisitions.
\textit{Human disagreement.} On CIFAR-10H~\cite{peterson2019}, with about 50 human labels per test image, the fuzzy model has the highest
correlation with human disagreement under label noise (Spearman 0.28 vs.\ 0.11 for MSP); on the clean networks it is above MSP but below
Gorbett and Blanchard (0.33 vs.\ 0.37 with ResNet-32 and 0.30 vs.\ 0.34 with VGG-16).

% =====================================================================
\section{Discussion and Limitations}
The features add most when the output is least trustworthy: under label noise, the final layers memorize wrong labels and the softmax
becomes overconfident, while the probes still reveal that the earlier blocks disagree. On easy problems (MNIST, SVHN), the softmax is
already near the ceiling and the internal reading adds nothing. On HAM10000, the binary output is already a strong signal once the
threshold is taken into account, so internal features mainly provide the explanation.
The detector is trained on the network's errors; with few errors, as in the 3,000-lesion subset, the variance across seeds is large.
The view features multiply the inference cost by $V$, although the single-pass variant retains most of the gain.
Finally, all detectors were trained by cross-validation on the test set of the same network; whether the rules transfer to another
architecture is left for future work.

% =====================================================================
\section{Conclusion}
Reading a frozen network from the inside, through the path of its decision across blocks, its stability across views and the familiarity
of the input, gives a misclassification score that is as accurate as the best competing detectors and, unlike them, can be stated as a
dozen weighted fuzzy rules. The gain over the network's own confidence is largest when that confidence is least reliable, as under label
noise and in skin-lesion classification. Future work will study whether the rules transfer across architectures, for instance with
interval type-2 memberships that represent the disagreement between networks, and validate the referral policy with dermatologists.

% Duplo-cego: sem agradecimentos na submissão.
% \section*{Acknowledgment}

\bibliographystyle{IEEEtran}
\bibliography{refs}

\end{document}
